

% ---------------------------------------------
\subsection{Errors for Stokes Parameters}
\label{sec: stokes_errors}

For the data calibrated in this project (\bfourNS, \bfiveNS, \bsevenNS), the errors for \SIns, \SQ and \SU were estimated by averaging the maximum standard deviation in four different regions far from the bright source at the centre. This average value was applied as a constant error at every pixel in the respective error maps. Table \ref{tab: source_info} gives the errors at each wavelength. The errors for the \bsix data were estimated using the same method. 

% ---------------------------------------------
\subsection{Polarization Fraction and Angle}

The following equations in this Chapter are from \cite{CoudeEquations}, unless specified otherwise. 

Using the debiased polarized intensity, the debiased polarization fraction can be calculated by:
\begin{equation}
    \POLFeq = \POLIeq/I,
    \label{eq: POLF}
\end{equation}
\noindent where $I$ is \SI intensity. The uncertainty on polarization fraction is:

\begin{equation}
    \POLFerrEQ = \POLFeq \left[\left(\frac{\POLIerrEQ}{\POLIeq}\right)^2 + \left(\frac{\SIerrEQ}{I}\right)^2\right]^{1/2}.
     \label{eq: POLF_error}
\end{equation}
   


The polarization position angle, $\theta$ is calculated as: 
\begin{equation}
\theta = \half \tan ^{-1} \frac{U}{Q},
    \label{eq: PA}
\end{equation}
\noindent where $\theta$ has units of radians, and an uncertainty calculated by:

\begin{equation}
    \PAerrEQ = \frac{1}{2} 
    \frac{\sqrt{(U\SQerrEQ)^2 + (Q\SUerrEQ)^2}}{Q^2 + U^2}.
    \label{eq: PA_error}
\end{equation}
 

\noindent When measuring the angles, the zero-degree position is defined as North, with the polarization angle increasing counterclockwise from North to East. The coordinate system was previously shown in see Figure \ref{fig: LinearPol}. 



% ---------------------------------------------







